\section{The independence architecture: Lemmas A--D}\label{sec:architecture}

Write the governor indicator as a functional of two layers of data:
\[
\mathbf 1_{\D}(m)\;=\;\Phi\big(\mathcal A(m),\,\mathcal R(m)\big),
\]
where $\mathcal A(m)$ is the \emph{anatomy} (the multiset of sizes
$u_i=\log p_i/\log m$ of the prime factors $p_i>m^{\delta}$, with
multiplicity) and $\mathcal R(m)$ is the \emph{digit data} (the residues
$m\bmod p_i^{j}$ at those primes, equivalently the fractional parts
$\fr{m/p_i^j}$). Membership in $\D$ is: every large prime power
$p^e\,\|\,m$ accrues $e$ carries; for the critical \emph{band} primes
$p=m^{u}$, $u\in(1/3,1/2)$, exactly one digit position ($j=2$) is
genuinely random --- position $j=1$ never carries beyond trivialities and
positions $j\ge3$ are deterministic at this size --- so a band prime is a
single near-fair coin. The proof of
Corollary~\ref{cor:reduction}'s correlation statement is a cascade of four
independence claims, stated here in the form used downstream.

\subsection{Lemma A: small primes are ignorable}

\begin{lemma}[A]\label{lem:A}
\statusKA\quad
For $0<\delta\le1$ write $n(n-1)=A_nB_n$ with
$P^+(A_n)\le x^{\delta}<P^-(B_n)$. Then
\[
\#\Big\{n\le x:\ v_p(A_n)>\ka_p(n)\ \text{for some }p\le x^\delta\Big\}
\;\ll\; x\,e^{-c/\delta},
\]
and the contribution of such $n$ to all correlation statements below is a
removable $\varepsilon$-loss as $\delta\to0$ (after $x\to\infty$).
Moreover, discarding $n$ with $p^2\mid n(n-1)$ for some $p>x^\delta$ costs
$\ll x^{1-\delta}$.
\end{lemma}

\begin{proof}[Source and the one-line change]
Ford--Konyagin's small-prime counting (\cite{FK21}, \S2) applies verbatim
on the $n$-side: for $p^v\mid n$ the low base-$p$ digits of $n$ are $0^v$
and failure of the carry budget forces an atypical digit pattern in the
remaining positions, with a binomial digit-count bound. The single
adaptation for the shifted side: for $q^v\mid n-1$ the low base-$q$ digits
of $n$ are $(1,0^{v-1})$ rather than $0^v$; these positions likewise
contribute no carries, and the identical counting applies to the remaining
positions. Union over the two sides, over $p\le x^\delta$ and over $v$, as
in \cite{FK21}, Proposition~1.
\end{proof}

\begin{remark}[a structural no-go worth recording]
The loss $e^{-c/\delta}$ is a \emph{fixed positive density} for fixed
$\delta$. Any proof design that requires a cofactor to be
$x^{\delta}$-smooth pays a Dickman-type probability
$\rho(c/\delta)$, which loses to $e^{-c/\delta}$ for every $\delta$: no
smooth-cofactor design can work. This dead end is recorded so it is not
re-entered.
\end{remark}

\subsection{Lemma B: anatomy independence}

\begin{lemma}[B]\label{lem:B}
\statusC\quad
For finite unions of anatomy boxes
$A,B\subseteq\{(u_1\ge u_2\ge\cdots):\sum_iu_i\le1\}$,
\[
\frac{1}{\log x}\sum_{n\le x}\frac{1}{n}\,
\mathbf 1\big[\mathcal A(n)\in A\big]\,
\mathbf 1\big[\mathcal A(n-1)\in B\big]
\;=\;\mathrm{PD}(A)\,\mathrm{PD}(B)+O\big((\log x)^{-c}\big)
\]
at almost all scales, where $\mathrm{PD}$ denotes the
Billingsley--Dickman measure of the box.
\end{lemma}

\noindent\emph{Citation route.} Tao--Ter\"av\"ainen
\cite{TT25}, building on Pilatte \cite{Pil23}, prove quantitative binary
correlations of $1$-bounded multiplicative functions with power-of-log
savings at almost all scales; their Theorem~1.7 is the prototype
(consecutive smooth pairs, $\rho(u)\rho(v)$). The hypotheses are
asymmetric: $g_1$ must be equidistributed (real case) or non-pretentious,
while $g_2$ need only be $1$-bounded multiplicative --- so the awkward
function goes in the $g_2$ slot.

\noindent\emph{The $z$-interpolation bridge.} Anatomy-box indicators are
not multiplicative, but $n\le x$ has at most $3$ prime factors above
$x^{1/3}$. For $z\in\C$ let $g_z$ be completely multiplicative with
$g_z(p)=z$ on a band and $1$ off it; then $g_z(n)=z^{B(n)}$ with
$B(n)\in\{0,1,2,3\}$ the band count, and Vandermonde interpolation at $4$
values of $z$ recovers each $\mathbf 1[B(n)=j]$ as a finite linear
combination of $1$-bounded multiplicative functions --- TT-admissible
objects. General band-anatomy indicators are finite
products/combinations across finitely many bands.

\noindent\emph{Pending fine print} (\S\ref{sec:remaining}, items B1--B3):
uniformity in the $x$-dependent band edges; the exact conclusion shape of
their Theorem~3.1; and the verbatim statement of that theorem (the
verification was interrupted and must be completed). Fallback for
Theorem~\ref{thm:lower}: Hildebrand's elementary stable-set method
\cite{Hil85} gives a positive-density lower bound for the anatomy layer
without TT machinery.

\subsection{Lemma C: within-side digit equidistribution, twisted}

\begin{lemma}[C]\label{lem:C}
\statusO{} (standard toolbox)\quad
Conditional on the anatomies of $n$ and $n-1$, the digit data
$\mathcal R(n)$ at the interior positions equidistributes at its product
rate, uniformly against any $1$-bounded multiplicative weight $g(n-1)$:
the Ford--Konyagin count of $n$ with prescribed large-prime data and digit
successes, weighted by $g(n-1)$, factors as
$(\text{FK main term})\times(\text{mean of }g)+o(1)$.
\end{lemma}

\noindent\emph{Two structural corrections baked into the statement.}
(i) Equidistribution is claimed for digit positions strictly below the
top; the \emph{top digit is deterministic} given the size ratio, following
the explicit piecewise law in $\fr{1/u}$ of \cite{FK21} ($p^2\in(n,2n]$:
forced success; $p^2\in(2n/3,n]$: forced failure; alternating at slot
boundaries). (ii) For two band primes the joint residues live on thin
progressions rather than filling the product grid; independence is
restored only on ensemble average over the prime pairs --- the phenomenon
treated by the Friedlander--Iwaniec analysis of simultaneous fractional
parts, the hardest ingredient here.

\noindent\emph{Where the new sums live.} Inserting $g(n-1)=g(pm-1)$ into
FK's Fourier detection turns the $m$-sums into
$\sum_{m\sim x/p}g(pm-1)\,\e{hm/p^{a-1}}$: mean values of bounded
multiplicative functions in progressions to moduli $p^{a-1}$, needed on
average over $p\sim x^u$ --- Bombieri--Vinogradov-for-multiplicative-functions
territory (Granville--Shao \cite{GS}, Matom\"aki--Radziwi\l\l--Tao
\cite{MRT}). Interface note: Lemma C is scheduled after the Lemma-D gate
since its final form depends on which version of D it must feed.

\subsection{Lemma D: cross-side digit independence (the hard core)}

\begin{lemma}[D]\label{lem:D}
\statusO\quad
Conditional on the anatomies, the digit layers of the two sides are
jointly independent: the $(n-1)$-side digit conditions (by
Lemma~\ref{lem:sandwich}, the carry conditions of $n-1$ at its own primes)
equidistribute after conditioning on the anatomy \emph{and} the digit
outcome of the $n$-side. Equivalently, $\sum_{n\le x}F_1(n)F_2(n-1)$
factors as the product of means when the $F_i$ carry anatomy and digit
data.
\end{lemma}

\noindent\emph{Reduction to arithmetic form.} Parametrize the dominant
configuration by band primes $p=x^u\mid n$, $q=x^{u'}\mid n-1$,
$u,u'\in(1/3,1/2)$, writing $n=pa$, $n-1=qb$. The binding $(n-1)$-side
digit condition at $q$ is
\[
n\equiv1\ (\mathrm{mod}\ q)
\quad\text{and}\quad
\frac{n-1}{q}\bmod q\in\Big[\frac q2,\,q\Big),
\]
a union of residue classes modulo $q^2$ lying above $1\bmod q$, with
$q^2\in[x^{2/3},x]$ --- beyond classical dispersion for general
coefficients. The harmonic dual of the mod-$q^2$ refinement is the family
of \emph{depth-one (Fermat-quotient) characters}
$\chi_\lambda(n)=\e{\lambda\fq(n)/q}$ developed in \S\ref{sec:dls}; the
homomorphism property $\fq(pa)=\fq(p)+\fq(a)$ makes all phases completely
bilinear. The quantitative core is:

\begin{estimate}[D$^\dagger$; the residual core, AP-native form]\label{est:Ddagger}
\statusO\quad
There is $c>0$ such that for each fixed $0<|\lambda|\le(\log x)^{C}$, on
average over band primes $q\sim Q=x^{u'}$,
\[
\sum_{m\le M_q}W_q(m)\,\e{\frac{\lambda m}{q}}\ \ll\ M_q\,(\log x)^{-c},
\qquad
W_q(m)=w(1+qm),\quad M_q=\frac{x-1}{q},
\]
where $w$ is the $1$-bounded, centered population-and-digit decoration of
the $n$-side (anatomy indicator times centered digit weight).
\end{estimate}

\noindent The program's decomposition of Estimate~D$^\dagger$, with
statuses:
\begin{center}
\begin{tabular}{ll}
\toprule
Layer & Status \\
\midrule
$|\lambda|\ge\Lambda_0$ (deep harmonics) &
\statusP{} --- Proposition~\ref{prop:largelambda} \\
digit layer, minor arcs & \statusP{} --- Theorem~\ref{thm:minor} \\
digit layer, major arcs (two families) &
counted; modulo Hypotheses~\ref{hyp:E1},~\ref{hyp:E2} \\
anatomy layer, generic & exact unwinding --- \S\ref{sec:anatomy} \\
anatomy layer, fine classes & \statusN{} --- \S\ref{sec:remaining} \\
\bottomrule
\end{tabular}
\end{center}

\noindent\emph{Hardness diagnosis (settled).} The statement is parity-free
(indicators are non-negative, positive-mean, local), so the Chowla parity
wall does not apply; the true hardness class is Bombieri--Vinogradov-type
equidistribution to large moduli on average. Orthogonality over the deep
character family is a Wieferich-type condition (\S\ref{sec:dls}),
pointwise hopeless, average-friendly. A direct citation of the
Tao--Ter\"av\"ainen/Pilatte machinery is blocked: multiplicativity enters
their argument at centring and dilation, and digit conditions are not
dilation-invariant; only Pilatte's $1$-bounded spectral engine survives as
a potential tool.

\subsection{Assembly and trims}\label{subsec:assembly}

\begin{proposition}[assembly]\label{prop:assembly}
\statusN\quad
Lemma~A (with $\delta\to0$ last), Lemma~B for the anatomy product,
Lemmas~C and D for the conditional digit products, and the
Ford--Konyagin computation of $c_1$ (\cite{FK21}, \S7) combine, with the
error bookkeeping of the cascade
\[
\Pr^{\log}\big[n\in\D,\ n-1\in D_0^{(\pm)}\big]
=\Pr[\mathcal A(n)\in\cdot]\;\Pr[\mathcal A(n-1)\in\cdot]
\;\Pr[\mathcal R(n)\,|\,\mathcal A]\;\Pr[\mathcal R(n-1)\,|\,\text{rest}]
+o(1),
\]
to give Theorem~\ref{thm:main}. The limit constant is Ford--Konyagin's
verbatim, squared.
\end{proposition}

\noindent\emph{Trim registry} (density losses accrued by the lower-bound
track; each is $O(\eta)$ and acceptable for Theorem~\ref{thm:lower}):
\begin{enumerate}
\item top-of-band trim: exclude $u'\in(1/2-\eta,1/2]$;
\item corner trim: exclude $u+u'\in(1-\eta'',1]$
(needed by Proposition~\ref{prop:digit-assembly});
\item margin costs: $\le6\times10^{-5}$ measured (Lemma~\ref{lem:margins});
\item a single positive-proportion anatomy subclass suffices for
Theorem~\ref{thm:lower}, and the automatic-carry band
$q^2\in(n,2n]$ (log-measure $0.288$ of slot-$2$ successes) is free.
\end{enumerate}
